\scanpage{66}
\mspage{53}
% archive sheet number: 65

la flèche $v_{U,i}$ est $\unsure{\in|\text{ds}} C$. Le diagramme $(\alpha)$ s'envoie dans le diagramme similaire pour $f':X'\mathbin{-}Y'$, $i'=u(i)$, et l'ouvert $U'\subset I'_{/i'}$ engendré par $u(U)$ :

\diagram{\begin{tikzcd}[ampersand replacement=\&,column sep=small,row sep=small]
 \& X_U\arrow[rr,dash]\arrow[dr,"{f_U}"']\arrow[ddd,"{c^{X'}_{\deleted{\unsure{X}}U}}"'] \& \& X(i)\arrow[dr,dash]\arrow[ddd,"{\mathrm{id}}" near end]\arrow[dl,phantom,"{\text{cocart}}" description] \& \& \\
 \& \& Y_U\arrow[rr]\arrow[ddd,"{c^{Y'}_{\deleted{\unsure{Y}}U}}"'] \& \& Y_{U,i}\arrow[dr,"{v_{U,i}}"]\arrow[ddd,"{c^{f'}_{U,i}}"'] \& \\
 (\alpha\alpha) \& \& \& \& \& Y(i)\arrow[ddd,"{\mathrm{id}}"]\\
 \& X'_{U'}\arrow[rr]\arrow[dr,"{f'_{U}}"'] \& \& X'(i')\arrow[dr,dash]\arrow[dl,phantom,"{\text{cocart}}" description] \& \& \\
 \& \& Y'_{U'}\arrow[rr] \& \& Y'_{U',i'}\arrow[dr,"{v_{U',i'}}"'] \& \\
 \& \& \& \& \& Y'(i')\ .
\end{tikzcd}}

Par cofinalité, \deleted{\unsure{les}} deux flèches verticales $c^{X'}_U$ et $c^{Y'}_U$ sont iso, \deleted{\unsure{donc}} plus de $X(i)\overset{\mathrm{id}}{\mathbin{-}} X'(i')$, donc \deleted{$\unsure{Y_{U,i}}$} la 4\textsuperscript{e} flèche verticale du cube, $c^{f'}_{U,i}$, est iso. Comme $v_{U',i'}$ est $\in C$, $v_{U,i}$ aussi, on \unsure{gagne|a gagné}.

\deleted{\unsure{Prouvons} Nous n'avons pas utilisé le th. 3.1, mais nous allons l'utiliser pour prouver b).}

\inserted{Prouvons b).} Soit donc \deleted{$f$ \unsure{une}}
\[
 f:X\longrightarrow Y
\]
une flèche dans $C_I$, et soit
\[
 f'=u_!(f):u_!(X)=X'\longrightarrow u_!(Y)=Y',
\]
prouvons que c'est dans $C_{I'^{\unsure{\circ|0}}}$, donc que
